\section{Limitations}
\label{sec:limitations}

We are explicit about what \wcabench{} cannot support, because a benchmark's
credibility depends on the honesty of its caveats.

\paragraph{Extrapolation risk (\task{4}).} Human-limit estimation has no
observable ground truth and is intrinsically vulnerable to structural change.
Historical data cannot anticipate a future jump driven by new hardware (for
example magnetic cubes), new training regimes, or rule changes. The
discrepancy between the three-by-three estimates of the change-point model
($3.02$~s), the Gaussian-process/extreme-value model ($3.29$~s) and exponential
decay ($3.74$~s), and the community anchor ($\approx 2.37$~s), should be read as
a warning, not as a resolution. Limits reported here are model-conditional
statements, not facts about human ability.

\paragraph{Regional imbalance.} Access to sanctioned competitions is uneven
across countries and continents. Because features such as competition count and
opponent strength are partly proxies for opportunity, cross-regional
generalization claims are weak. We report continental strata for transparency,
but small strata (for example Oceania, South America) contain too few test
instances for reliable conclusions, and several stratified cells are
undefined.

\paragraph{Self-selection and confounding (\task{5}).} Participation in an
event is a choice correlated with ability. Correlation-based transfer
estimates therefore overstate effects, and even the DID proxy relies on a
parallel-trends assumption that self-selection may violate. The
``identifiable pair count'' is a measure of how much the data support an
estimate, not a guarantee of a causal effect. \wcabench{} should not be used to
make causal claims about training without an explicit identification argument.

\paragraph{Temporal non-stationarity.} The training window (2003--2022) differs
from the test window (2025--2026) in hardware, rules and participant
demographics. Models may degrade in ways that reflect genuine distribution
shift rather than modelling error; the falling \task{2} Kendall $\kendall$
across test slices ($0.90 \to 0.64$) is consistent with exactly this.

\paragraph{Sampled small-mode results.} The headlines in
Section~\ref{sec:results} come from a sampled test window run on a single
machine with one GPU (Section~\ref{sec:compute}). Absolute values will move
when the full rolling-window protocol is executed on all $1{,}719{,}290$ test
records, and ranking ties (such as the \task{2} $0.7673$ tie) may resolve
differently; the three-seed table (Section~\ref{sec:results-multiseed}) is
included precisely to expose this sampling variance but does not substitute for
the full run. We report the sampled numbers because they are real and
verifiable, not because they are final; the released pipeline supports the full
run.

\paragraph{Imbalanced, sparse classes.} \DNF{} prediction and skill-transfer
estimation both suffer from severe imbalance and sparsity. Even conditioning on
rounds that reach the decision point, the positive rate varies from single
digits in some events to about $95\%$ in three-by-three blindfolded, and many
event pairs have too few common competitors for stable estimation.

\paragraph{Native backends and compute reproducibility.} This release runs the
tree baselines on native XGBoost and the GNN on its native CUDA path
(Section~\ref{sec:compute}); every report records the device, wall-clock time
and CPU/GPU hours actually used. Readers reproducing a baseline should hold the
library versions and device fixed: a different backend (for example a
scikit-learn fallback for the tree models, or a CPU-only graph kernel) will not
reproduce these numbers exactly, and the earlier fallback-backed results
differed materially from the native ones.

\paragraph{Single-snapshot data.} Even with a fixed export (v2.0.2), the \wca{}
occasionally corrects historical results, and the database grows continuously:
each weekly export carries a new sequential version label and additional
competitions. We pin the 2026-09-21 snapshot and record it, so the exact scale
figures (for example $6{,}909{,}454$ results) are snapshot-specific and a newer
export will show slightly larger numbers. Readers comparing against the current
official download should treat our counts as a frozen baseline rather than a
live value.

\section{Ethics and Broader Impact}
\label{sec:ethics}

\paragraph{Permitted uses.} \wcabench{} is intended for scientific research and
teaching: sports data analysis, benchmark methodology, and reproducible
evaluation of prediction and inference methods. It may also serve the
speedcubing community directly, for example to help competitors understand
their own progression or organisers to reason about round structure.

\paragraph{Prohibited uses.} The accompanying data card forbids, in the
strongest terms, the use of these data or models for:
\begin{itemize}
  \item \textbf{gambling, betting or wagering} of any kind, including
        predictive services marketed to bettors;
  \item \textbf{discriminatory screening, profiling or shaming} of individual
        competitors;
  \item \textbf{impersonating} the \wca{} or other official bodies, or
        fabricating competition results; and
  \item redistributing individual-level derived data without anonymization.
\end{itemize}
These prohibitions are not decorative. Aggregate competition records are
public, but they are records of identifiable people, and their aggregation
enables inference that individuals did not consent to. The risk that a
predictive benchmark is repurposed for betting is real enough that we treat its
prevention as a design requirement.

\paragraph{Attribution.} All derived publications must carry the \wca{}
attribution notice: “This information is based on competition results owned and
maintained by the World Cube Association, published at
\url{https://worldcubeassociation.org/results}.” Code is released under
Apache-2.0; data rights remain with the \wca{}.

\paragraph{Privacy and opt-out.} Although the underlying data are public,
individuals may have legitimate reasons to withdraw. We therefore commit to an
opt-out mechanism: a competitor may contact the maintainers to request removal
of their individual-level derived features from released artifacts. Aggregate
statistics---which do not single out an individual---are retained, so the
benchmark's scientific claims remain reproducible while the individual's
choice is respected. We also recommend that challenge organisers restate the
prohibited uses in their rules and that the \wca{} Results Team remain a
standing point of contact.

\paragraph{Fairness considerations.} Because regional opportunity is unequal,
a model trained predominantly on data from well-represented regions may perform
worse for under-represented ones, and this could, if misused, disadvantage
competitors from those regions. We mitigate by stratifying by continent and by
refusing to publish individual-level rankings as products. We do not claim that
stratification resolves the underlying inequity; it only makes it visible.

\paragraph{Broader impact.} We hope \wcabench{} nudges sports machine learning
towards more rigorous evaluation: time-aware splits, frozen statistics,
stratified reporting, effect sizes and explicit identification assumptions.
The main societal risk is misuse for betting or individual targeting, which the
data card and the opt-out mechanism are designed to discourage.
